% Source: https://artofproblemsolving.com/community/c3h2163492_fun_math_problems
% Stable identity: aops-topic-2163492-q01
Solving for $y$: the first line is $y=-2x+2$, with slope $-2$, and the second is $y=\dfrac{1}{2}x$, with slope $\dfrac{1}{2}$. The slopes are different, so the lines are not parallel, and their product is $-2\cdot\dfrac{1}{2}=-1$, so the lines are $\boxed{\text{perpendicular}}$. They meet at $\left(\dfrac{4}{5},\dfrac{2}{5}\right)$.

\begin{asy}
usepackage("amsmath");
size(5cm);
defaultpen(fontsize(10pt));
import graph;
xaxis(-1.5,2.5,Arrow(TeXHead)); yaxis(-1.2,3.2,Arrow(TeXHead));
draw((-0.5,3)--(1.6,-1.2),blue+linewidth(0.9));
draw((-1.5,-0.75)--(2.4,1.2),red+linewidth(0.9));
pair X=(0.8,0.4);
draw(X+0.25*unit((1,-2))--X+0.25*unit((1,-2))+0.25*unit((2,1))--X+0.25*unit((2,1)));
dot(X);
label("$2x+y=2$",(1.6,-1.2),E,blue); label("$x-2y=0$",(2.4,1.2),N,red);
label("\mbox{}",(min(currentpicture,user=true)+max(currentpicture,user=true))/2,black); // keeps the text after the figure black when LaTeX typesets the labels inline
\end{asy}
